\chapter{Considerações finais}
\label{cap:consideracoes-finais}

Este trabalho apresentou a concepção, a implementação e a validação inicial do BoraFut, uma solução móvel para apoiar a organização do futebol amador em quadras públicas do Distrito Federal. O desenvolvimento percorreu levantamento e priorização de requisitos, prototipação, identidade visual, arquitetura cliente-servidor, modelagem de dados, implementação do aplicativo e do backend, automação de testes e implantação em ambiente piloto.


\section{Alcance dos objetivos}
O objetivo geral de desenvolver e avaliar uma solução multiplataforma foi alcançado dentro do recorte proposto. O produto resultante reúne descoberta de quadras, criação e compartilhamento de futs, confirmação e auditoria de presença, fila, formação e rotação de times, condução da partida em tempo real, avaliações e mecanismos de intervenção para situações imprevistas. O registro estruturado dos eventos permite reconstruir a operação e produzir indicadores sem depender apenas da memória dos participantes.

Os objetivos específicos declarados na Seção~\ref{sec:objetivos} foram cumpridos em graus distintos. O projeto seguiu práticas estabelecidas de Engenharia de Software, da abordagem híbrida de processo à modelagem arquitetural em níveis e ao registro das decisões de projeto. A instrumentação de eventos de domínio no backend foi implementada e mostrou seu valor de forma concreta: foram os próprios registros de chegada e de partida que permitiram apurar os indicadores do teste de campo, sem recorrer à memória dos presentes. A qualidade de uso foi trabalhada no design, orientado por heurísticas consolidadas, e verificada junto a usuários por meio da SUS. O catálogo georreferenciado de quadras públicas foi construído com atributos qualificadores e aberto à avaliação dos jogadores, ainda que restrito à região administrativa do piloto. O MVP foi disponibilizado em ambiente piloto, com aplicativo distribuído nas faixas de teste das lojas sobre infraestrutura de produção. O processo e as decisões foram documentados nesta monografia e nos apêndices, incluindo especificações de casos de uso, dicionário de dados e modelo lógico. Por fim, a avaliação dos resultados e a indicação de melhorias constituem o Capítulo~\ref{cap:validacao-resultados} e o restante deste capítulo.

Dois objetivos merecem ressalva quanto ao alcance. O mapeamento de quadras cobriu a região do piloto, e não o conjunto do Distrito Federal, o que era esperado no recorte adotado mas limita a generalização da cobertura. A avaliação de usabilidade, por sua vez, incidiu sobre o protótipo, e não sobre a aplicação final: o teste de campo observou a operação do produto implementado, porém sem aplicar novamente um instrumento padronizado de usabilidade. Medir a usabilidade percebida do aplicativo em uso continuado permanece, portanto, como etapa a cumprir.

A avaliação ocorreu em duas etapas. Cinco sessões realizadas em maio com o protótipo de alta fidelidade produziram média SUS bruta de 85,5; uma resposta reconhecidamente interpretada de forma invertida motivou análise de sensibilidade, cuja média foi 94,4 entre os quatro formulários restantes. Em agosto, o primeiro teste em produção reuniu 27 chegadas e 15 partidas efetivamente iniciadas. O aplicativo não se tornou gargalo da atividade, e as contingências observadas puderam ser tratadas pelas funcionalidades disponíveis.


\section{Limitações}
Esses resultados devem ser interpretados com cautela. As amostras foram pequenas e de conveniência; o teste de usabilidade avaliou um protótipo sob mediação; o teste de campo ocorreu em uma única quadra, com participação ativa dos autores na operação e sem grupo de controle. Não houve acompanhamento longitudinal suficiente para medir retenção, recorrência dos organizadores ou impacto sobre a ocupação das quadras. A ausência de bug bloqueador observado em uma noite tampouco demonstra ausência de defeitos no sistema.

O teste de campo evidenciou que o principal risco imediato é comportamental. Parte dos jogadores utilizou o aplicativo na preparação ou na chegada, mas não manteve interação durante o fut. Como o telefone não acompanha todos os participantes dentro da quadra, a solução precisa continuar funcional com poucos aparelhos e, ao mesmo tempo, tornar claros os momentos em que alguma ação é necessária. Melhorias no encerramento da partida, na comunicação da rotação e no início da partida seguinte são prioritárias para reduzir dúvidas sem criar o papel de um mesário permanente.


\section{Trabalhos futuros}
Como trabalhos futuros, destacam-se a evolução das estatísticas dos jogadores, o refinamento contínuo dos temas claro e escuro, a ampliação da observabilidade em produção e a validação em futs recorrentes operados sem a presença dos autores. Somam-se a esses os itens que este trabalho deixou deliberadamente para depois: a expansão do catálogo de quadras para outras regiões administrativas do Distrito Federal, hoje restrito à Asa Norte; a filtragem dos resultados do mapa pela área visível, adiada por não ser necessária no volume atual de quadras; os requisitos registrados como posteriores ao produto mínimo no Project Model Canvas, entre eles campeonatos, gamificação e comunicação interna entre participantes; e as histórias de usuário marcadas como fora do escopo no Apêndice~\ref{apendice:us}, como a troca de endereço de e-mail, o resumo dos dados pessoais do titular, a verificação por mensagem de texto e o convite por e-mail ao participante cadastrado presencialmente.

Conclui-se que o BoraFut alcançou maturidade suficiente para apoiar um fut real em ambiente de produção e oferece uma base técnica consistente para evolução. A contribuição deste trabalho não se limita ao aplicativo: ela inclui a sistematização das regras informais do futebol amador, a incorporação dessas regras em um sistema auditável e a documentação de um processo de engenharia orientado por evidências. O passo seguinte é converter a aceitação inicial em uso recorrente, preservando a simplicidade necessária ao contexto social e físico das quadras públicas.
